%!TEX root = ../main.tex

\section{생성 모델 훈련}

데이터 세트가 로드되고 데이터 스키마가 정의되면 SDK는 모델 교육을 진행합니다.
이 단계에서 SDK는 데이터세트 크기와 사용 가능한 메모리를 기반으로 배치 크기와 학습률을 자동으로 선택하고 학습률 스케줄러를 사용하여 수렴을 더욱 최적화합니다.
사용자는 CPU와 GPU 하드웨어를 모두 활용하여 훈련을 가속화할 수 있으며, 훈련 실행을 투명하게 추적하고 관리할 수 있도록 실시간 진행 상황 모니터링이 제공됩니다.

다음에서는 먼저 핵심 모델 접근 방식에 대한 개요를 제공한 다음 지원되는 LLM에 대한 세부 정보를 제공합니다. 마지막으로 SDK에 구현된 개인정보 보호 메커니즘이 제시됩니다.

\subsection{자동 배열 테이블 형식 개요}

생성기 아티팩트는 ~\cite{tabularargn}에 자세히 설명된 TabularARGN(Tabular Auto-Regressive Generative Network)을 사용하여 훈련됩니다.
해당 아키텍처는 이산화된 기능에서 작동하고 범주형 교차 엔트로피 손실을 사용하여 최적화되는 얕은 임의 차수 자동 회귀 모델을 제공합니다.
기본적으로 TabularARGN은 생성 프로세스를 자동 회귀 프레임워크에 매핑하여 데이터의 결합 확률 분포를 더 간단한 조건부 확률의 산물로 추정합니다.
다중 테이블 데이터 세트를 처리하기 위해 모델은 자동 회귀 조건부 확률 프레임워크를 확장하여 데이터 스키마에 정의된 관련 테이블 간의 종속성을 캡처합니다.
이러한 시나리오에서 SDK는 테이블이 처리되는 순서를 명시적으로 정의하여 테이블 간의 상위-하위 관계를 설정합니다.
이런 방식으로 상위 테이블의 레코드는 \emph{context} 임베딩으로 압축된 다음 하위 테이블의 조건부 생성에 사용됩니다.

TabularARGN의 중요한 측면은 테이블 내의 열이 모델링되는 순서입니다.
SDK에서 \texttt{mostly.train()}를 호출할 때 사용자는 훈련을 위한 두 가지 열 순서 전략, 즉 고정 순서와 유연한 순서(또는 임의 순서) 중에서 선택할 수 있습니다.
기본적으로 유연한 옵션이 활성화되어 있습니다. 이 모드에서 SDK는 각 훈련 배치에 대한 열 순서를 무작위로 섞습니다. 이 접근 방식을 통해 모델은 기능의 하위 집합에 따라 조건이 지정된 데이터를 생성하는 방법을 학습할 수 있으므로 데이터 대치, 조건부 시뮬레이션 또는 공정성 조정과 같은 작업에 다용도로 사용할 수 있습니다.
유연한 교육을 통해 더 폭넓은 생성 기능을 사용할 수 있지만 학습 작업도 더욱 복잡해집니다. 이러한 추가적인 복잡성으로 인해 결과 합성 데이터의 정확성이 저하될 수 있습니다. 고정 열 순서를 선택하면 학습이 단순화되고 경우에 따라 정확도가 향상될 수 있습니다. 그러나 일단 고정 순서 모드로 훈련되면 모델은 임의의 기능 하위 집합에서 유연하게 데이터를 생성하는 기능을 잃게 됩니다.
이는 생성 시 더 높은 정확성과 더 많은 유연성 사이의 절충안을 제시하며, 의도한 애플리케이션에 따라 최적의 선택이 가능합니다.
% In the SDK, when calling \texttt{mostly.train}, users can control the auto-regressive order of the training process through the \verb|enable_flexible_generation| flag.
% If set to \verb|False|, the order of the columns is kept fixed. 
% When set to \verb|True|, the SDK dynamically shuffles the order of the columns for each training batch. 
% This allows the model to learn conditional probabilities based on any subset of features.



\subsection{LLM 절단 지원}

HuggingFace 플랫폼에서 호스팅되는 LLM 미세 조정을 지원합니다.
% As mentioned in Section~\ref{subsec:data_encoding}, the user can choose to encode columns using the language encoding types. 
% When doing this, our SDK provides the option of processing language encoded columns with LLMs, as long as the user has a HuggingFace access token with read permissions. 
섹션~2에서 설명한 대로 텍스트 열 모델링을 위해 사전 훈련된 LLM을 활용하도록 선택할 수 있습니다. 이 경우 사용자는 처음부터 훈련되는 기본 LSTM 모델에 의존하는 대신 데이터를 미세 조정하기 위해 원하는 LLM을 선택합니다. 이 옵션은 사용자에게 읽기 권한이 있는 HuggingFace 액세스 토큰이 있는 한 SDK에서 사용할 수 있습니다.
현재 HuggingFace에서 사용 가능한 모든 텍스트 생성 LLM을 지원합니다.

\subsection{개인 정보 보호}

TabularARGN 훈련 절차는 훈련 데이터의 개별 기록을 보호하기 위한 강력한 개인 정보 보호 장치로 설계되었습니다. 민감한 정보를 기억하게 될 수 있는 과적합의 위험을 최소화하기 위해 모델은 전용 개인 정보 보호 중심 전략을 사용합니다.
\begin{itemize}
\item \textbf{조기 중지:} 훈련 중 데이터 세트의 최대 10\%를 검증 세트로 따로 둡니다. 각 epoch가 끝날 때 검증 손실을 계산합니다. 검증 손실이 연속 $N$ epoch 동안 개선되지 않으면(기본값: 5) 훈련을 조기에 중단합니다. 관측된 검증 손실이 가장 낮았던 모델 가중치를 최종 모델로 유지합니다.
\item \textbf{드롭아웃 정규화:} 각 네트워크 계층에는 25\%의 드롭아웃 비율이 포함되어, 훈련 중 뉴런의 일부를 무작위로 비활성화함으로써 과적합 위험을 줄입니다.
\end{itemize}

\paragraph{차등 개인 정보 보호}
SDK는 DP 보장을 통해 합성 데이터 생성기를 교육하는 옵션을 제공합니다. DP가 활성화되면 훈련 절차는 원본 데이터 세트의 개별 레코드가 개인 정보 보호 예산 매개변수 $\epsilon$로 정량화되는 것처럼 훈련된 모델에 수학적으로 제한된 영향을 미치도록 보장합니다. 개인 정보 보호 예산은 훈련 전반에 걸쳐 모니터링되며, 이 임계값에 도달하면 자동으로 훈련을 중단하도록 상한 $\epsilon$ 제한을 설정할 수 있습니다.

SDK의 차등 개인 정보 보호는 그라데이션 클리핑, 노이즈 추가, 무작위 데이터 로딩 및 개인 정보 계정을 통해 개인 정보 보호를 보장하는 Opacus 라이브러리를 사용하여 구현됩니다. 사용자는 요구 사항에 따라 개인 정보 보호와 모델 유틸리티의 균형을 맞추기 위해 노이즈 수준과 클리핑 매개변수를 구성할 수 있습니다.

비DP 훈련은 효율성과 정확성을 극대화하기 위한 기본값이지만, DP 옵션은 모델 성능과 계산 효율성에 영향을 미칠 수 있다는 점을 이해하고 엄격한 수학적 개인 정보 보호를 요구하는 사용자에게 제공됩니다.
